% ===================================================================
%  TABLEAU N° 6 — La Maison intérieure. — Les Meubles. — La
%  Nourriture. — L'Éclairage.
%  Feuillets 37 a 44 du fac-simile = folios 35 a 42.
%  Correspondance : folio imprime = numero de feuillet - 2.
%
%  Les cles « %%K » sont les MEMES que dans texto/io : c'est par elles
%  que la page de lecture met les deux textes en regard. Le francais
%  decoupe ce tableau en CINQ sections numerotees en chiffres romains
%  (I. La Chambre a coucher — II. Le Salon — III. La Salle a Manger —
%  IV. La Cuisine — V. La Salle de Bain), la ou l'ido en fait NEUF
%  (« La Dormo-chambro », « La laborkorbo... », « La Salono », « La
%  Bubacho », « La Manjo-chambro », « La Dineo-servo », « La
%  Koqueyo », « La Vazaro », « Altra objekti... », « La Balneyo »).
%  L'auteur remet la numerado des alineas a 1 au debut de chaque
%  section francaise, comme il le fait pour chaque chambro ido : on
%  klavize ca que videsas, sen korektar; outils/ceni.py appliquera
%  l'indice de scene apres coup.
%
%  ATTENTION — LE FAC-SIMILE FRANCAIS EST UNE PRISE DE VUE, non un
%  passage au scanner : l'echelle change d'une page a l'autre. La
%  calibration metrique de kalibro-fr.tex est donc une MEDIANE.
%
%  Coquille conservee (§7 de la consigne) : au feuillet 40 (folio 40),
%  alinea 4 de « La Cuisine », l'imprime coupe « fourneau de cui- »
%  puis « suine » — la composition ajoute un u parasite a « cuisine ».
%  Verifie a l'agrandissement (2600 px) : la lettre est bien un u.
%  Coquille conservee egalement au titre III : « LA SALLE A MANGER, »
%  se termine par une virgule et non un point dans le fac-similé.
% ===================================================================

% --- Feuillet 37 (folio 35, ouverture de tableau) -------------------
\begin{VUpage}[37]{35}
%%K t06-apar-1 apar
\VUsaut{6.0mm}
\VUtitre{15.0pt}{60}{TABLEAU N\textsuperscript{o} 6}
\VUsaut{1.4mm}
\VUtitre{11.4pt}{0}{La Maison intérieure. --- Les Meubles.}
\VUsaut{0.4mm}
\VUtitre{11.4pt}{0}{La Nourriture. --- L'Éclairage.}
\VUsaut{2.2mm}

%%K t06-apar-2 apar
\VUblancAlinea
La maison est terminée. L'oncle de Jean est installé\nl
dans son nouveau domicile, dont nous connaissons la\nl
distribution. Voyons un peu comment il a meublé sa\nl
maison. Nous visitons d'abord :

%%K t06-tit-1 sub
\VUpk{10.2pt}{40}{I. --- La Chambre à coucher.}
\VUsaut{3.0mm}

%%K t06-c1-01-1 p
\VUblancAlinea
1. --- Nous arrivons mal à propos, car la bonne est\nl
occupée à faire le \VUgras{ménage.}

%%K t06-c1-02-1 p
\VUblancAlinea
2. --- Sur la \VUgras{table de toilette} \textsuperscript{(1)} sont éparpillés\nl
les différents objets avec lesquels on se lave, on se\nl
peigne, en un mot, on fait sa toilette. Ce sont la\nl
\VUgras{cuvette} \textsuperscript{(2)} et le \VUgras{pot à eau} \textsuperscript{(3)}, la \VUgras{brosse à che}\cc
\VUgras{veux} \textsuperscript{(4)}, le \VUgras{peigne} \textsuperscript{(5)}, le \VUgras{savon} \textsuperscript{(6)} et l'\VUgras{éponge} \textsuperscript{(7)};\nl
enfin un petit \VUgras{flacon} \textsuperscript{(8)} pour l'eau dentifrice.

%%K t06-c1-03-1 p
\VUblancAlinea
3. --- Par terre, devant la toilette, voici le \VUgras{seau} \textsuperscript{(9)}\nl
pour recevoir les eaux sales.

%%K t06-c1-04-1 p
\VUblancAlinea
4. --- Nous apercevons dans l'\VUgras{antichambre} \textsuperscript{(10)}\nl
quelques \VUgras{portemanteaux} \textsuperscript{(13)} et deux \VUgras{parapluies} \textsuperscript{(11)}\nl
dans le \VUgras{porte-parapluies} \textsuperscript{(12)}.

%%K t06-c1-05-1 p
\VUblancAlinea
5. --- De l'autre côté de la porte se trouve\nl
l'\VUgras{armoire à glace} \textsuperscript{(14)}, qui contient le \VUgras{linge} \textsuperscript{(15)}.

%%K t06-c1-06-1 p
\VUblancAlinea
6. --- Profitons de ce que la \VUgras{femme de chambre} \textsuperscript{(16)}\nl
fait le \VUgras{lit} \textsuperscript{(17)} pour en étudier les différentes parties.\nl
Le \VUgras{bois de lit} \textsuperscript{(20)} contient un bon \VUgras{sommier} \textsuperscript{(21)} sur\nl
lequel Justine va déposer un \VUgras{matelas} \textsuperscript{(22)} de laine.\nl
Elle prendra ensuite sur les chaises les deux \VUgras{draps} \textsuperscript{(23)}\nl
et la \VUgras{couverture} \textsuperscript{(24)}. Puis elle placera à la tête du
\parplein
\end{VUpage}

% --- Feuillet 38 (folio 36) -----------------------------------------
\begin{VUpage}[38]{36}
%%K t06-c1-06-1 p suite
\VUcontinue
\VUgras{lit} le \VUgras{traversin} \textsuperscript{(25)} et l'\VUgras{oreiller} \textsuperscript{(26)} qui sont avec\nl
l'\VUgras{édredon} \textsuperscript{(27)} sur la \VUgras{descente de lit} \textsuperscript{(32)}. Un coup\nl
de plumeau pour secouer la poussière des \VUgras{rideaux} \textsuperscript{(19)}\nl
et du \VUgras{ciel de lit} \textsuperscript{(18)}, et voilà une partie de la besogne\nl
terminée.

%%K t06-c1-07-1 p
\VUblancAlinea
7. --- Il faudra alors mettre un peu d'ordre sur la\nl
\VUgras{table de nuit} \textsuperscript{(28)}, remonter le \VUgras{réveille-matin} \textsuperscript{(29)},\nl
préparer la \VUgras{veilleuse} \textsuperscript{(30)} et emporter la \VUgras{bougie} \textsuperscript{(31)}\nl
pour nettoyer le bougeoir.

%%K t06-c1-08-1 p
\VUblancAlinea
8. --- La \VUgras{corbeille à ouvrage} \textsuperscript{(34)}, auprès du\nl
\VUgras{tabouret} \textsuperscript{(33)}, a aussi grand besoin d'être rangée. Il\nl
faudra plier la \VUgras{broderie} \textsuperscript{(35)}, serrer dans la \VUgras{table à}\nl
\VUgras{ouvrage} \textsuperscript{(38)} le \VUgras{dé}, le \VUgras{fil} \textsuperscript{(36)}, les \VUgras{aiguilles} \textsuperscript{(37)} et les\nl
\VUgras{ciseaux} \textsuperscript{(39)} et fermer avec la \VUgras{clef} \textsuperscript{(40)} le couvercle\nl
de la table.

%%K t06-c1-09-1 p
\VUblancAlinea
9. --- Sur le \VUgras{guéridon} \textsuperscript{(41)} du milieu, Justine renou\cc
vellera l'eau de la \VUgras{carafe} \textsuperscript{(42)}. Elle mettra du sucre\nl
dans le \VUgras{sucrier} \textsuperscript{(43)}, nettoiera la \VUgras{pince à sucre} \textsuperscript{(44)}\nl
et enlèvera la \VUgras{brosse à habits} \textsuperscript{(46)}.

%%K t06-c1-10-1 p
\VUblancAlinea
10. --- Le côté droit de la chambre lui donnera\nl
moins de travail, car la \VUgras{chaise longue} \textsuperscript{(47)} et son\nl
\VUgras{coussin} \textsuperscript{(48)} sont à leur place; et, comme il ne fait\nl
pas froid, rien n'est dérangé parmi les accessoires de\nl
la \VUgras{cheminée} \textsuperscript{(49)}. \VUgras{Pelle} \textsuperscript{(50)}, \VUgras{pincettes} \textsuperscript{(51)}, \VUgras{che}\cc
\VUgras{nets} \textsuperscript{(55)}, \VUgras{garde-cendre} \textsuperscript{(56)} sont propres; l'\VUgras{écran} \textsuperscript{(54)}\nl
n'a pas été dérangé et le \VUgras{coffre à bois} \textsuperscript{(52)} est bien\nl
garni de \VUgras{bois} \textsuperscript{(53)}.

%%K t06-c1-11-1 p
\VUblancAlinea
11. --- Il faudra cependant épousseter la \VUgras{pendule} \textsuperscript{(57)},\nl
les \VUgras{candélabres} \textsuperscript{(58)} et les \VUgras{vide-poches} \textsuperscript{(59)}, puis\nl
essuyer la \VUgras{glace} \textsuperscript{(60)}.

%%K t06-c1-12-1 p
\VUblancAlinea
12. --- Quant au \VUgras{berceau} \textsuperscript{(61)} du bébé, il est déjà\nl
tout prêt.

%%K t06-c1-12-2 p
\VUblancAlinea
\textit{Nota.} --- \textit{Le professeur pourra prendre les mots les}\nl
\textit{plus importants de ce vocabulaire pour décrire une}\nl
\textit{chambre d'hôtel infiniment plus simple.}
\end{VUpage}

% --- Feuillet 39 (folio 37) -----------------------------------------
\begin{VUpage}[39]{37}
%%K t06-tit-2 sub
\VUpk{10.2pt}{40}{II. --- Le Salon.}
\VUsaut{3.0mm}

%%K t06-c2-01-1 p
\VUblancAlinea
1. --- Voici maintenant le salon. C'est une très jolie\nl
\VUgras{pièce.} La cheminée, qui se trouve dans l'angle de\nl
droite, est ornée d'une fine \VUgras{statuette} \textsuperscript{(1)}, de deux\nl
beaux \VUgras{vases} \textsuperscript{(3)} et drapée avec un \VUgras{lambrequin} \textsuperscript{(2)} de\nl
velours.

%%K t06-c2-02-1 p
\VUblancAlinea
2. --- Pour éviter que les étincelles du foyer ne\nl
mettent le feu au tapis, on a placé devant l'âtre un\nl
\VUgras{pare-étincelles} \textsuperscript{(4)} en cuivre.

%%K t06-c2-03-1 p
\VUblancAlinea
3. --- Un élégant \VUgras{bureau} \textsuperscript{(5)} de dame est ouvert à\nl
côté. C'est là que la \VUgras{maîtresse de maison} \textsuperscript{(23)} vient\nl
écrire ses lettres.

%%K t06-c2-04-1 p
\VUblancAlinea
4. --- La fenêtre est garnie de \VUgras{rideaux de}\nl
\VUgras{vitrage} \textsuperscript{(6)} relevés dans des \VUgras{embrasses} \textsuperscript{(8)} accro\cc
chées aux \VUgras{patères} \textsuperscript{(7)} et de grands rideaux d'étoffe\nl
surmontés d'une \VUgras{draperie} \textsuperscript{(9)}.

%%K t06-c2-05-1 p
\VUblancAlinea
5. --- Dans l'embrasure de la fenêtre, un \VUgras{cache}\cc
\VUgras{pot} \textsuperscript{(11)} contient une \VUgras{plante} \textsuperscript{(10)} verte, un palmier\nl
magnifique dont les feuilles s'étendent presque jusqu'à\nl
la \VUgras{lampe} \textsuperscript{(12)} à pétrole.

%%K t06-c2-06-1 p
\VUblancAlinea
6. --- L'\VUgras{abat-jour} \textsuperscript{(13)} de cette lampe a été fait par\nl
Marie, la charmante \VUgras{jeune fille} \textsuperscript{(14)} qui est au\nl
\VUgras{piano} \textsuperscript{(15)}. Assise bien droite sur le \VUgras{tabouret} \textsuperscript{(16)},\nl
elle étudie un morceau. Elle est fort adroite et aussi\nl
très soigneuse, comme le prouve son \VUgras{casier à}\nl
\VUgras{musique} \textsuperscript{(17)}, qui est parfaitement rangé.

%%K t06-c2-07-1 p
\VUblancAlinea
7. --- Son frère Paul cherche un \VUgras{livre} \textsuperscript{(19)} dans la\nl
\VUgras{bibliothèque} \textsuperscript{(18)}. C'est un \VUgras{jeune garçon} \textsuperscript{(20)} re\cc
muant et insupportable. En s'accrochant à la biblio\cc
thèque pour atteindre le livre qui est placé trop haut,\nl
il risque de faire tomber le \VUgras{buste} \textsuperscript{(21)} qui se trouve\nl
au-dessus.

%%K t06-c2-08-1 p
\VUblancAlinea
8. --- Il profite, pour faire cette sottise, de ce que sa\nl
maman est occupée à causer avec une amie, une\nl
\VUgras{dame} \textsuperscript{(22)} qui est venue passer quelques jours chez
\parplein
\end{VUpage}

% --- Feuillet 40 (folio 38) -----------------------------------------
\begin{VUpage}[40]{38}
%%K t06-c2-08-1 p suite
\VUcontinue
elle. La maman est assise sur le \VUgras{canapé} \textsuperscript{(24)}, à côté\nl
du \VUgras{fauteuil} \textsuperscript{(25)}, et son amie est sur le \VUgras{siège} \textsuperscript{(27)}, dont\nl
on ne voit que le \VUgras{dossier} \textsuperscript{(26)}.

%%K t06-c2-09-1 p
\VUblancAlinea
9. --- Le grand \VUgras{cadre} \textsuperscript{(30)} suspendu au mur con\cc
tient le \VUgras{portrait} \textsuperscript{(29)} d'une parente. Dans l'\VUgras{album} \textsuperscript{(32)}\nl
posé sur la table sont réunies les \VUgras{photographies} \textsuperscript{(33)}\nl
des autres membres de la famille. Les enfants l'ont\nl
feuilleté tout à l'heure, car les \VUgras{chaises} \textsuperscript{(31)} sont en\cc
core groupées autour de la table.

%%K t06-c2-10-1 p
\VUblancAlinea
10. --- La \VUgras{potiche} \textsuperscript{(34)} en porcelaine qui se trouve\nl
près du \VUgras{coupe-papier} \textsuperscript{(35)} a été donnée à Marie pour\nl
sa fête, et le \VUgras{paravent} \textsuperscript{(36)} qui garnit l'angle de la\nl
pièce a été rapporté de Chine par son oncle.

%%K t06-c2-11-1 p
\VUblancAlinea
11. --- Égayé par les jolis \VUgras{tableaux} \textsuperscript{(37)} accrochés\nl
à la \VUgras{cloison} \textsuperscript{(42)}, éclairé par le \VUgras{lustre} \textsuperscript{(38)} du plafond\nl
dont les \VUgras{lampes électriques} \textsuperscript{(39)} jettent, le soir, une\nl
si joyeuse clarté, ce salon paraît fort agréable. Le\nl
\VUgras{tapis} \textsuperscript{(40)} moelleux étendu sur le parquet et la \VUgras{son}\cc
\VUgras{nette électrique} \textsuperscript{(41)} si pratique, achèvent de le\nl
rendre très confortable.

%%K t06-tit-3 sub
\VUsaut{3.0mm}
\VUpk{10.2pt}{40}{III. --- La Salle à Manger,}
\VUsaut{3.0mm}

%%K t06-c3-01-1 p
\VUblancAlinea
1. --- L'heure du dîner est arrivée. Toute la famille,\nl
à l'exception de Marie, est rassemblée autour de la\nl
table. La salle à manger est doucement chauffée par\nl
un excellent petit \VUgras{poêle} \textsuperscript{(1)} en faïence, muni d'un\nl
mince \VUgras{tuyau} \textsuperscript{(2)}.

%%K t06-c3-02-1 p
\VUblancAlinea
2. --- Cette salle ne comprend pas beaucoup de\nl
meubles. Le long du mur est suspendue, au-dessus\nl
de l'\VUgras{escabeau} \textsuperscript{(4)}, une petite \VUgras{fontaine} \textsuperscript{(3)} décorative.\nl
Tout auprès se trouve le \VUgras{dressoir} \textsuperscript{(5)} (la desserte),\nl
sur lequel on dépose les plats froids, les desserts et\nl
les différents ustensiles de table dont on ne se sert\nl
pas momentanément.

%%K t06-c3-03-1 p
\VUblancAlinea
3. --- C'est ainsi que nous voyons, sur la tablette\nl
supérieure, un \VUgras{poulet rôti} \textsuperscript{(7)}, un \VUgras{poisson} \textsuperscript{(8)} et une
\parplein
\end{VUpage}

% --- Feuillet 41 (folio 39) -----------------------------------------
\begin{VUpage}[41]{39}
%%K t06-c3-03-1 p suite
\VUcontinue
\VUgras{coupe} \textsuperscript{(10)} contenant des \VUgras{fruits} \textsuperscript{(9)}. Au-dessous, voici\nl
le \VUgras{fromage} \textsuperscript{(11)}, le \VUgras{pain} \textsuperscript{(12)}, l'\VUgras{huilier} \textsuperscript{(13)}, qui com\cc
prend tout ce qu'il faut pour assaisonner la \VUgras{salade} :\nl
l'\VUgras{huile}, le \VUgras{vinaigre}, le \VUgras{sel} \textsuperscript{(14)} et le \VUgras{poivre} \textsuperscript{(15)}.\nl
Enfin, sur la tablette inférieure, se trouve, à côté du\nl
\VUgras{moutardier} \textsuperscript{(16)}, un \VUgras{gâteau} \textsuperscript{(17)}.

%%K t06-c3-04-1 p
\VUblancAlinea
4. --- La pendule de la salle à manger, que l'on\nl
nomme \VUgras{cartel} \textsuperscript{(20)}, a une forme très particulière. Elle\nl
est enchâssée dans un cadre en bois, qui contient aussi\nl
le \VUgras{baromètre} \textsuperscript{(19)} et le \VUgras{thermomètre} \textsuperscript{(18)}.

%%K t06-c3-05-1 p
\VUblancAlinea
5. --- Tout autour de la pièce court une jolie\nl
\VUgras{boiserie} \textsuperscript{(21)} en chêne ciré. Une \VUgras{portière} \textsuperscript{(22)} en\nl
étoffe ferme suffisamment la porte qui conduit à la\nl
véranda. C'est là que la famille ira prendre le café\nl
après le dîner. Déjà les \VUgras{tasses} \textsuperscript{(24)} et leurs \VUgras{sou}\cc
\VUgras{coupes} \textsuperscript{(25)} sont préparées sur le \VUgras{buffet} \textsuperscript{(23)}.

%%K t06-c3-06-1 p
\VUblancAlinea
6. --- Au-dessus de la table est fixée au plafond\nl
une \VUgras{suspension} \textsuperscript{(26)} dont la puissante lampe éclaire\nl
le soir toute la salle à manger.

%%K t06-c3-07-1 p
\VUblancAlinea
7. --- Occupons-nous maintenant de la \VUgras{table} \textsuperscript{(27)} elle\cc
même. Quand la famille est entrée, le \VUgras{couvert} était\nl
mis. Sur la \VUgras{nappe} \textsuperscript{(28)} étaient posés, à la place de\nl
chaque convive, une \VUgras{assiette} \textsuperscript{(33)}, une \VUgras{serviette} \textsuperscript{(29)},\nl
une \VUgras{cuillère} \textsuperscript{(34)}, une \VUgras{fourchette} \textsuperscript{(35)}, un \VUgras{cou}\cc
\VUgras{teau} \textsuperscript{(36)} et un \VUgras{verre} \textsuperscript{(32)}. Il y avait aussi des \VUgras{bou}\cc
\VUgras{teilles} \textsuperscript{(30)} pour le vin et une \VUgras{carafe} \textsuperscript{(31)} pour l'eau.

%%K t06-c3-08-1 p
\VUblancAlinea
8. --- Le \VUgras{domestique} \textsuperscript{(39)} a tout d'abord apporté la\nl
\VUgras{soupière} \textsuperscript{(37)}, dans laquelle un peu de \VUgras{potage} (soupe),\nl
fume encore. Il sert maintenant le premier \VUgras{mets} \textsuperscript{(38)},\nl
un plat de viande. Il présentera ensuite ceux que\nl
nous avons déjà nommés, puis, après le dessert, il\nl
profitera de ce que ses maîtres seront passés dans la\nl
véranda pour desservir la table et mettre en ordre la\nl
salle à manger.

%%K t06-c3-08-2 p
\VUblancAlinea
\textit{Nota.} --- \textit{Les termes usuels se rapportant à la nourriture}\nl
\textit{sont donnés ici d'une façon très sommaire. La liste en sera}\nl
\textit{complétée dans les deux tableaux intitulés : l'Hôtel (n\textsuperscript{o} 13)}\nl
\textit{et le Marché (n\textsuperscript{o} 15).}
\end{VUpage}

% --- Feuillet 42 (folio 40) -----------------------------------------
\begin{VUpage}[42]{40}
%%K t06-tit-4 sub
\VUpk{10.2pt}{40}{IV. --- La Cuisine.}
\VUsaut{3.0mm}

%%K t06-c4-01-1 p
\VUblancAlinea
1. --- Dans la cuisine, où nous jetons un coup d'œil\nl
par la porte entr'ouverte, nous apercevons un pitto\cc
resque désordre. Devant l'\VUgras{âtre} \textsuperscript{(1)} où le \VUgras{feu} \textsuperscript{(3)} pétille,\nl
le \VUgras{tourne-broche} \textsuperscript{(5)} est installé. Un beau \VUgras{rôti} \textsuperscript{(7)}\nl
est à la \VUgras{broche} \textsuperscript{(6)} et achève de cuire.

%%K t06-c4-02-1 p
\VUblancAlinea
2. --- Le \VUgras{soufflet} \textsuperscript{(8)} et la \VUgras{pelle à charbon} \textsuperscript{(9)},\nl
dont on vient de se servir, sont restés par terre ainsi\nl
que les \VUgras{bûches} \textsuperscript{(10)}.

%%K t06-c4-03-1 p
\VUblancAlinea
3. --- Le dessus de la cheminée est en ordre;\nl
la \VUgras{lanterne} \textsuperscript{(11)}, le \VUgras{porte-allumettes} \textsuperscript{(13)}, dans lequel\nl
sont les \VUgras{allumettes} \textsuperscript{(12)}, le \VUgras{chandelier} \textsuperscript{(14)} et le \VUgras{bou}\cc
\VUgras{geoir} \textsuperscript{(15)}, contenant une \VUgras{chandelle} \textsuperscript{(16)}, sont à leur\nl
place. Mais le grand \VUgras{seau à charbon} \textsuperscript{(18)}, plein de\nl
\VUgras{charbon de terre} \textsuperscript{(17)}, que la \VUgras{cuisinière} \textsuperscript{(23)} vient\nl
d'aller remplir à la cave, est resté au milieu de la\nl
pièce.

%%K t06-c4-04-1 p
\VUblancAlinea
4. --- Il est vrai que la pauvre femme est obligée\nl
de se hâter pour que son déjeuner soit prêt à temps.\nl
Pour le moment, elle est près du \VUgras{fourneau de cui}\cc
\VUgras{suine} \textsuperscript{(19)} et tourne une \VUgras{sauce} \textsuperscript{(24)} qu'il ne faut pas\nl
qu'elle laisse brûler :

%%K t06-c4-05-1 p
\VUblancAlinea
5. --- Dans le \VUgras{four} \textsuperscript{(20)} cuit un gâteau, qu'elle a\nl
préparé pour le goûter des enfants. Au-dessus du\nl
fourneau sont suspendues la \VUgras{cuillère à pot} \textsuperscript{(21)} et\nl
l'\VUgras{écumoire} \textsuperscript{(22)}.

%%K t06-c4-06-1 p
\VUblancAlinea
6. --- La \VUgras{batterie de cuisine} \textsuperscript{(25)}, accrochée dans\nl
le fond de la pièce, est très propre. Les belles \VUgras{casse}\cc
\VUgras{roles} \textsuperscript{(26)} de cuivre, le \VUgras{pot au lait} \textsuperscript{(27)}, la \VUgras{passoire} \textsuperscript{(28)}\nl
et la \VUgras{râpe} \textsuperscript{(29)} ont été nettoyés avec soin.

%%K t06-c4-07-1 p
\VUblancAlinea
7. --- La \VUgras{boîte au sel} \textsuperscript{(30)} est posée sur le petit\nl
buffet à côté de la \VUgras{théière} \textsuperscript{(31)} et de la \VUgras{bonbonne} \textsuperscript{(32)}\nl
à huile. Le \VUgras{tiroir} \textsuperscript{(33)} contient probablement les cou\cc
verts et les couteaux de cuisine.
\end{VUpage}

% --- Feuillet 43 (folio 41) -----------------------------------------
\begin{VUpage}[43]{41}
%%K t06-c4-08-1 p
\VUblancAlinea
8. --- Le \VUgras{balai} \textsuperscript{(34)} et le \VUgras{plumeau} \textsuperscript{(35)} sont dans\nl
un coin. La cuisinière vient de se servir du \VUgras{billot} \textsuperscript{(36)};\nl
elle l'a laissé près de la fenêtre.

%%K t06-c4-09-1 p
\VUblancAlinea
9. --- Un \VUgras{essuie-mains} \textsuperscript{(37)} est suspendu au mur,\nl
à côté de l'\VUgras{entonnoir} \textsuperscript{(38)}. Dans cette partie de la\nl
cuisine sont relégués le \VUgras{gril} \textsuperscript{(39)}, le \VUgras{hachoir} \textsuperscript{(40)} et la\nl
\VUgras{planche à hacher} \textsuperscript{(41)}. Puis, au-dessus, sur une\nl
\VUgras{planche} \textsuperscript{(42)}, sont des \VUgras{pots} \textsuperscript{(43)}, le \VUgras{pot au feu} \textsuperscript{(44)} avec\nl
son \VUgras{couvercle} \textsuperscript{(45)}, et la \VUgras{marmite} \textsuperscript{(46)}. La \VUgras{poêle} \textsuperscript{(47)}\nl
est suspendue à côté.

%%K t06-c4-10-1 p
\VUblancAlinea
10. --- Seul le \VUgras{robinet} \textsuperscript{(48)} de cuivre jette une note\nl
brillante dans ce coin. L'\VUgras{évier} \textsuperscript{(49)}, sur lequel est\nl
placée la grosse \VUgras{cruche} \textsuperscript{(50)}, doit être bien commode\nl
avec son petit placard où l'on renferme le \VUgras{ton}\cc
\VUgras{neau} \textsuperscript{(51)} à vinaigre muni de sa \VUgras{cannelle} \textsuperscript{(52)}, la \VUgras{boîte}\nl
\VUgras{à ordures} \textsuperscript{(53)} et la \VUgras{vaisselle} \textsuperscript{(54)} sale.

%%K t06-c4-11-1 p
\VUblancAlinea
11 --- Le grand \VUgras{cuvier} \textsuperscript{(55)} dans lequel on lave les\nl
\VUgras{légumes} \textsuperscript{(58)} et le \VUgras{chaudron} \textsuperscript{(56)} de fonte posé sur\nl
le \VUgras{carrelage} \textsuperscript{(57)} doivent eux aussi trouver place\nl
dans ce placard.

%%K t06-c4-12-1 p
\VUblancAlinea
12. --- La cuisinière ne manque certes pas de four\cc
neaux, car voilà encore, sur le coin de la table, un\nl
\VUgras{fourneau à gaz} \textsuperscript{(59)} sur lequel chauffe de l'eau con\cc
tenue dans une petite casserole. La \VUgras{vapeur} \textsuperscript{(60)} qui\nl
s'en échappe nous indique qu'elle doit bouillir. Il est\nl
probable que cette eau va servir à faire le café, car\nl
voici à l'autre bout de la table la \VUgras{cafetière} \textsuperscript{(64)} et le\nl
\VUgras{moulin à café} \textsuperscript{(63)}.

%%K t06-c4-13-1 p
\VUblancAlinea
13. --- Le fourneau est relié au \VUgras{bec de gaz} \textsuperscript{(62)}\nl
par un long \VUgras{tuyau de caoutchouc} \textsuperscript{(61)}.

%%K t06-c4-14-1 p
\VUblancAlinea
14. --- La \VUgras{femme de ménage} \textsuperscript{(66)}, qui vient aider\nl
la cuisinière, prépare les légumes pour le dîner.\nl
Lorsqu'ils seront épluchés et lavés, elle les mettra\nl
dans la \VUgras{bassine} \textsuperscript{(65)} en attendant l'heure de les faire\nl
cuire.
\end{VUpage}

% --- Feuillet 44 (folio 42) -----------------------------------------
\begin{VUpage}[44]{42}
%%K t06-tit-5 sub
\VUpk{10.2pt}{40}{V. --- La Salle de Bain.}
\VUsaut{3.0mm}

%%K t06-c5-01-1 p
\VUblancAlinea
1. --- Il ne nous reste plus qu'une seule indiscrétion\nl
à commettre pour connaître les principales pièces de\nl
l'appartement : voir l'installation de la \VUgras{salle de}\nl
\VUgras{bain.} Ce ne sera pas long, car elle n'est pas grande,\nl
et nous saurons vite que la \VUgras{baignoire} \textsuperscript{(1)} a un\nl
bon \VUgras{chauffe-bain} \textsuperscript{(2)}, que le \VUgras{porte-savon} \textsuperscript{(3)} est\nl
à portée de la main de la personne qui se baigne\nl
et que le \VUgras{porte-serviette} \textsuperscript{(5)} doit permettre de faire\nl
sécher très facilement les \VUgras{serviettes de toilette} \textsuperscript{(4)}.

%%K t06-c5-02-1 p
\VUblancAlinea
2. --- Cette salle a ses cloisons recouvertes de\nl
\VUgras{faïence} \textsuperscript{(6)}, ce qui a permis d'installer un \VUgras{appareil}\nl
\VUgras{à douche} \textsuperscript{(7)} sans craindre que l'eau, qui rejaillit en\nl
s'en servant, ne détériore les murs. Un petit \VUgras{bassin} \textsuperscript{(8)}\nl
en zinc, qui sert à prendre les bains de pieds, complète\nl
l'ameublement.

\VUsaut{6.0mm}
\VUfilet{22mm}
\end{VUpage}
